% ---- pg-587 （书页 575） ----

户的联系, 设计更好的引擎并建立适当的设备, 创建更好的医疗程序, 改进过程控制, 并且发展有效的信用评分法则.

在传统的统计应用中, 样本容量相对较小, 数据是小心地搜集到的, 样本分析结果为统计推断提供了基础, 虽然也处理异常情况但不是首要关注, 通常高度依赖模型. 在数据挖掘中, 样本容量巨大, 数据是分散的, 有历史关系的(日复一日地记录), 样本用于训练、验证和检验(非正式的推断), 关注异常情况, 通常不建立模型. 不仅如此, 数据的准备(包括数据的搜集、评估和整理, 变量的定义和选择)通常是费力的任务, 大约在数据挖掘的工作中占 60\%$\sim$80\% 的比例.

数据挖掘的问题大致可分为以下几个内容:

\noindent$\bullet$\quad 分类(离散结果):

\noindent\hspace*{1.5em}谁可能改用另一家移动通信服务公司?

\noindent$\bullet$\quad 预测(连续结果):

\noindent\hspace*{1.5em}这栋房子的适当估价是多少?

\noindent$\bullet$\quad 关联/市场篮子分析:

\noindent\hspace*{1.5em}脱脂奶是否经常与低脂奶酪一起被购买?

\noindent$\bullet$\quad 聚类:

\noindent\hspace*{1.5em}是否存在有相同购买习惯的顾客?

\noindent$\bullet$\quad 描述:

\noindent\hspace*{1.5em}星期四, 杂货店的顾客通常同时购买玉米片和软饮料.

鉴于数据挖掘问题的特性, 不难理解为什么本书中讨论的许多统计方法构成完备的数据挖掘软件的重要部分. 具体地说, 回归、判别和分类程序(线性法则、逻辑斯蒂回归、决策树等)以及聚类算法是非常重要的数据挖掘工具. 其他工具包括关联法则, 多重适应回归样条(MARS), $K$ 近邻算法, 神经网络, 遗传算法以及可视化等. 这些方法超出了本书的范围.\textsuperscript{①}

\mvsub{数据挖掘的过程}

数据挖掘的过程有一系列步骤. 这些步骤构成一种策略, 与任何其他的建模过程的策略没有本质的不同. 具体而言, 数据挖掘必须

\noindent 1.\quad 廓清问题, 确定目标.

\noindent 2.\quad 收集和准备合适的数据.

\noindent 3.\quad 从数据中发现可能的关联、未察觉的特性以及显而易见的异常情况, 以得到深入的认识.

\noindent 4.\quad 清理数据并对变量作适当变换.

\noindent 5.\quad 将数据划分为训练、验证和检验数据集.

\noindent 6.\quad 对训练数据建立模型.

\noindent 7.\quad 用验证数据评估模型的好坏, 必要时修正模型.

\noindent 8.\quad 基于验证数据或检验数据评估模型的表现. 将模型输出与原始对象作比较. 考察模型是否有用.

\noindent 9.\quad 应用模型.

\vspace{6pt}
\noindent\rule{\linewidth}{0.4pt}
\par\vspace{2pt}
{\zihao{6}\noindent ①\quad 要得到更多的关于数据挖掘的一般信息和数据挖掘工具, 见本章末尾的参考文献.\par}
